% ===== 原文 pp.9–11 =====
% [原文 p.9]

在此设定下，人们考虑如下形式的、具有非正则奇点（irregular singularity）的稳定可积联络
\[
\mathcal{A}=\mathrm{d}+T_n\frac{\mathrm{d}z}{z^{n}}+\cdots+T_1\frac{\mathrm{d}z}{z}
\]
（其中 $n>1$），以及如上一节所述的稳定抛物 Higgs 丛 $(E,\Phi)$——其 Higgs 场带有极点部分，例如
\begin{equation}
T_n\frac{\mathrm{d}z}{z^{n}}+\cdots+T_1\frac{\mathrm{d}z}{z}, \tag{2.17}
\end{equation}
并且其定义使得 $E$ 在 $p_j$ 处的纤维被留数 $\mathrm{res}(\Phi,p_j)$ 所保持；这里 $T_1,T_2,\dots,T_n$ 是群\footnote{通常假定矩阵 $T_n$ 可对角化且具有互不相同的特征值；这就排除了幂零（nilpotence）情形，而幂零情形会排除出现高阶正则奇点的可能，同时也使我们能够定义形式对角化规范变换，把联络化为由 (1.15) 给出的非正则型（irregular type）。}的李代数的元素。在考虑低秩 Higgs 丛时，我们有时采用 [23] 的记号，把复化联络（complexified connection）写成
\begin{equation}
\mathcal{A}=\mathrm{d}Q+T_1\frac{\mathrm{d}z}{z}, \tag{2.18}
\end{equation}
其中 $Q$ 是亚纯联络构成的对角方阵\footnote{在定义 $Q$ 之前，通常先施行形式对角化步骤。}；改变 (2.18) 中的 $Q$ 等价于改变 (2.17) 中的 $(T_n,T_{n-1},\dots,T_2)$。

\begin{defn}{2.19} $\mathcal{A}$ 的形式单值性（formal monodromy）为 $M_0:=\ee^{2\pi\ii T_1}$。此外，我们说 $T_1$ 是形式单值性的指数（exponent）。\end{defn}

% [Tutor Note: p.9 首现「形式单值性」（formal monodromy）：$M_0=\ee^{2\pi\ii T_1}$ 只由对角矩阵 $T_1$（即形式单值性的指数）决定；它与 (2.23) 的 Stokes 矩阵合起来，才构成极点处的局部单值性数据（对照 p.11 末与 p.12 开头对完整单值性数据的说明）。]

仿照 Simpson 在 [96] 中对抛物 Higgs 丛的构造，自然可以假定这些联络与 Higgs 场在全纯规范等价意义下可化为极点部分为对角形的对象。与正则情形类似，这些野 Higgs 丛的模空间可以作为超凯勒商得到 [24,26]。在本讲义中，我们沿用 [25] 的记号：型为 $(m,r_1,r_2,\dots)$ 的野 Higgs 丛是指带 $m$ 个极点、且各极点阶数为 $k_i:=r_i+1$ 的 Higgs 丛。为了理解 Higgs 丛的模空间，下文我们把注意力限制在 $\PP^{1}$ 上带极点的秩 2 Higgs 丛。

\begin{rem}{2.20} 值得考察极点数（计重数）恰为 4 的 Higgs 丛，因为它们与 Painlevé 方程的解密切相关：(A) 型 $(1,3)$ 给出 Painlevé II——见问题集 I 的 (5).iii.；(B) 型 $(2,1,1)$ 给出 Painlevé III——见例 2.4.1；(C) 型 $(4,0,0,0,0)$ 给出 Painlevé VI——见第 2.3 节）。\end{rem}

% [Tutor Note: p.9 注 2.20 中的 "see Example 2.4.1" 在本文编号体系里找不到对应条目（本文的例只用 2.x 这类编号），疑为对下文小节 2.4.1（p.11 "2.4.1. Higgs bundles with multiple poles."）的误引或笔误；正文照原样译作「例 2.4.1」。]

% [Tutor Note: p.9 注 2.20 末尾 "see Section 2.3)." 多出一个右括号，原文如此（500 dpi 页图复核）；译文照原样保留末尾的「）」。]

对于秩 2 的 Higgs 丛，由 [25, 注 9.12] 可知：当 $(m,r_1,r_2,\dots)$ 依次为 $(4,0,0,0,0)$（即注 2.20 中的 (C)）、$(3,1,0,0)$、$(2,1,1)$（即注 2.20 中的 (B)）、$(2,2,0)$ 以及 $(1,3)$（即注 2.20 中的 (A)）时，所得模空间的复维数为 2。为了说明野 Higgs 丛研究中的主要要素，本章余下部分我们将考虑一些具体的例子。读者若欲了解细节与完整证明，可参阅前面提到的 Boalch 的论文，以及 [4, 附录 A] 中出现的更多与弦论（string theory）相关的例子。最后，关于这一领域的未决问题，[6] 给出了一份（并不详尽的）把野 Higgs 丛与奇异几何联系起来的问题清单；而关于与 Painlevé 方程的关系，例如可参阅 [98] 及其中的参考文献。

% [原文 p.10]

\para{具有一个极点的 Higgs 丛}
下面我们通过在 $\PP^{1}$ 上固定点 $p_1=0$ 处带一个二阶极点的秩 2 Higgs 丛，给出 [23] 中引入的主要要素的构造（另见 [5, 附录 A]）。在此设定下，(2.18) 中的联络为
\begin{equation}
\mathcal{A}=\mathrm{d}z\left(\frac{T_2}{z^{2}}+\frac{T_1}{z}\right), \tag{2.21}
\end{equation}
其中 $T_2=\frac{\mathrm{d}Q}{\mathrm{d}z}z^{2}$，$T_1$ 是常对角矩阵。特别地，我们写
\[
Q=\begin{pmatrix}q_1(z)&0\\ 0&q_2(z)\end{pmatrix},
\]
并定义 $q_{ij}(z)$ 为 $q_i(z)-q_j(z)$ 的首项。当 $q_1(z)-q_2(z)=a/z+b$ 时，有
\[
q_{12}(z)=a/z,\qquad q_{21}(z)=-a/z.
\]

% [Tutor Note: p.10 "the number of anti-Stokes directions in any half-period is 1 = Mult(d_1)+…+Mult(d_l)"：200/500 dpi 页图与文字层均显示为数字「1」；但按上文（半周期是 $l$ 元组、按重数加权计数）此数应为 $l$。正文照原样排作 $1$。]

反 Stokes 方向（anti-Stokes direction）的集合 $\mathcal{A}\subset S^{1}$ 由这样的方向 $d\in S^{1}$ 组成：在由 $d$ 指定的射线上，或者 $q_{12}(z)\in\R_{<0}$，或者 $q_{21}(z)\in\R_{<0}$。这些方向正是当 $z$ 趋于 $0$ 时 $e^{(a/z+b)}$（相应地 $e^{-(a/z+b)}$）衰减最快的方向；并且，由于若 $\pm\frac{a}{z}\in\R_{<0}$ 就有 $\mp\frac{a}{z}e^{-i\pi}\in\R_{<0}$，故 $\#\mathcal{A}:=r=2$。注意：若对某方向 $d$ 有 $q_{12}(z)\in\R_{<0}$，则该方向上 $q_{21}(z)\notin\R_{<0}$（反之亦然）。我们称一个方向为半周期（half-period）：取 $l:=r/2$，半周期是由 $l$ 个连续的反 Stokes 方向构成的 $l$ 元组 $d=(d_1,\dots,d_l)\subset\mathcal{A}$。$d$ 的重数（multiplicity）$\mathrm{Mult}(d)$ 是支撑 $d$ 的根的个数。按重数加权计算时，任一（半）周期中反 Stokes 方向的个数为 $1=\mathrm{Mult}(d_1)+\cdots+\mathrm{Mult}(d_l)$。在 (2.21) 的设定下，令 $a=\tilde{r}\ee^{i\tilde{\theta}}$，则只有两个半周期：

\begin{itemize}[leftmargin=1.4em,itemsep=1pt,topsep=1pt]
\item 由 $q_{12}$ 得到方向 $d_1:=\ee^{i(\tilde{\theta}-\pi)}$：取 $z=r'\ee^{i(\tilde{\theta}-\pi)}$，就有 $q_{12}(z)=-\frac{\tilde{r}}{r'}\in\R_{<0}$；
\item 由 $q_{21}$ 得到方向 $d_2:=\ee^{-i\pi}d_1=\ee^{i\tilde{\theta}}$：取 $z=r'\ee^{i\tilde{\theta}}$，便得到 $q_{21}(z)=-\frac{\tilde{r}}{r'}\in\R_{<0}$。
\end{itemize}

对反 Stokes 方向 $d$，$d$ 的根（roots）是如下定义的集合：
\[
\mathrm{Roots}(d):=\{(ij)\mid q_{ij}(z)\in\R_{<0}\},
\]
其中「$q_{ij}(z)\in\R_{<0}$」是指沿方向 $d$ 而言。注意，Stokes 方向 $d_1$ 与 $d_2$ 分别定义了序 $q_1<_{d_1}q_2$ 与 $q_2<_{d_2}q_1$。此外，这些方向定义了 Stokes 扇区（Stokes sector）$\mathrm{Sect}_i:=(d_i,d_{i+1})$，其中下标 $i$ 模 $\ell$。对由 (2.21) 给出的 Higgs 丛，其根为 $\mathrm{Roots}(d_1)=\{(12)\}$ 与 $\mathrm{Roots}(d_2)=\{(21)\}$。

% [Tutor Note: p.10 原文把同一量既写作 $l$（"for $l:=r/2$"、"$l$-tuple"、"$d_l$"），又写作 $\ell$（"$i\bmod\ell$"）；本译本照原文分别保留 $l$ 与 $\ell$，读者宜按同一量理解。]

% [原文 p.11]

\begin{rem}{2.22} 对由 (2.18) 给出、在零点处有一个 3 阶极点的 Higgs 丛，写
\[
q_1-q_2=\frac{a}{z^{4}}+\frac{b}{z^{3}}+\frac{c}{z^{2}}+\frac{d}{z}+e,
\]
便得到 $q_{12}(z)=\frac{a}{z^{3}}$ 与 $q_{21}(z)=-\frac{a}{z^{3}}$。该系统有 6 个反 Stokes 方向；我们将在问题集中回到它。\end{rem}

% [Tutor Note: p.11 注 2.22 的数据彼此不协调：由 p.10 的 (2.21) 例（2 阶极点时 $q_{12}=a/z$，恰有 2 个反 Stokes 方向）可知 $\#\mathcal{A}=2d$，其中 $d$ 是 $q_{12}=a/z^{d}$ 的极点阶数；于是「6 个反 Stokes 方向」要求 $d=3$，即 $q_{12}=a/z^{3}$（与原文一致），故 $q_1-q_2$ 的首项按定义应为 $a/z^{3}$，而非原文所写的 $a/z^{4}$。疑原文首项级数 $a/z^{4}+\cdots$ 多写了一阶（或 "order 3" 实指 $n=4$）。正文照原文排。]

与方向 $d$ 相联系的 Stokes 因子（Stokes factor）构成的群为
\[
\mathrm{Sto}_d(\mathcal{A}):=\{K\in G\mid (K)_{ij}=\delta_{ij}\ \ \forall\,(ij)\notin\mathrm{Roots}(d)\},
\]
对秩 2 的 Higgs 丛，它是 $G=\GL(2,\C)$ 的 1 维幂幺子群。对于我们所涉及的所有方向 $d_1$ 与 $d_2$，相应的 Stokes 矩阵群\footnote{应当注意，Boalch 等作者（例如见 [2]）把半个周期数量的 Stokes 因子之乘积称为「Stokes 矩阵」，即 (2.23) 中的对象——在秩 2 的情形，这两种定义一致。}是 1 维子群
\begin{equation}
\mathrm{Sto}_{d_1}(\mathcal{A})=\left\{\begin{pmatrix}1&\tilde{z}\\ 0&1\end{pmatrix}\ \middle|\ \tilde{z}\in\C\right\},
\qquad
\mathrm{Sto}_{d_2}(\mathcal{A})=\left\{\begin{pmatrix}1&0\\ \tilde{z}&1\end{pmatrix}\ \middle|\ \tilde{z}\in\C\right\}. \tag{2.23}
\end{equation}

为研究野 Higgs 丛的模空间，需要在每个 Stokes 扇区中考虑基本解，再用 Stokes 矩阵把它们延拓出去；这些 Stokes 矩阵可以看作典型基本解之间的转移矩阵。不过我们不打算深入这一点，而建议读者参阅前面提到的 Boalch 的论文（例如见 [20,23--26]）。

\para{2.4.1.\ 多极点的 Higgs 丛}
在考虑除子 $D$ 上带有不同极点的 Higgs 丛时，为构造其模空间需要考虑更多的数据。下面我们通过研究带四个极点的秩 2 Higgs 丛，提及几个要素，让读者对这套理论有所体会。这里取除子 $D=\{p_1,\dots,p_4\}$，取 $V$ 为 $\PP^{1}$ 上的同构平凡（homomorphically trivial）丛，并令 $\nabla=\mathrm{d}-\mathcal{A}$ 为在 $D$ 上有极点的亚纯联络，其中
\[
\mathcal{A}=\sum_{i=1}^{4}\left(T^{i}_{k_i}\frac{\mathrm{d}z}{(z-p_i)^{k_i}}+\cdots+T^{i}_{1}\frac{\mathrm{d}z}{(z-p_i)^{1}}\right),
\]
这里 $T^{i}_{j}$ 是对角 $2\times2$ 矩阵。联络 $\nabla$ 在 $p_i$ 处有 $k_i$ 阶极点；若假定 $T^{1}_{1}+\cdots+T^{4}_{1}=0$，则它没有其他极点。为处理这类 Higgs 丛，我们对每个 $i$ 取 $\PP^{1}$ 上以 $p_i$ 为心、两两不交的开圆盘 $D_i$，以及 $D_i$ 上在 $p_i$ 处为零的坐标 $z_i$。于是，前面针对单个极点描述的局部理论在每个这样的圆盘上重复：对应于不同极点的 Stokes 矩阵之间由联络矩阵（connection matrix）相联系，而联络矩阵是用局部基本解定义的转移矩阵。这些联络矩阵与所谓的触手（tentacle）相关联。

% [Tutor Note: p.11 原文 "let $\nabla=\mathrm{d}-\mathcal{A}$ be a meromorphic connections …" 中 connections 为复数形式，应为单数 connection（原文笔误）；译文按单数处理。]

为完整描述带多个极点的非正则联络的单值性数据，需要每个极点处的 Stokes 矩阵与形式单值性，以及联系不同极点邻域内基本解的联络矩阵——为此还要选取极点之间的点与路径。
